% from ega4-16-fr.tex
\oldpage[IV-4]{46}
  tout homomorphisme $u$ de $\mathscr{P}_{X/S}^n(\mathscr{F})$ dans
  $\mathscr{G}$ correspondant biunivoquement à la famille de ses composés
  \[
    u_\mu:\mathscr{P}_{X/S}^n(\mathscr{F})
    \longrightarrow\mathscr{G}\longrightarrow\mathscr{G}_\mu.
  \]
  On a donc
  \[
    \operatorname{Diff}_{X/S}^n(\mathscr{F},\mathscr{G})
    =\prod_{\mu\in M}
    \operatorname{Diff}_{X/S}^n(\mathscr{F},\mathscr{G}_\mu),
  \]
  et par suite aussi
  \[
    \mathscr{D}\!iff_{X/S}^n(\mathscr{F},\mathscr{G})
    =\prod_{\mu\in M}
    \mathscr{D}\!iff_{X/S}^n(\mathscr{F},\mathscr{G}_\mu).
  \]
  \item[(ii)] Jusqu'ici, on n'a guère rencontré d'opérateurs différentiels
  $\mathscr{F}\to\mathscr{G}$ que lorsque $\mathscr{F}$ et $\mathscr{G}$ sont
  des $\mathcal{O}_X$-Modules localement libres de rang fini, auquel cas, leur
  structure se ramène localement, en vertu de (i), à celle du faisceau
  $\mathscr{D}\!iff_{X/S}$; cette dernière va être étudiée plus loin (16.11)
  dans un cas particulier.
\end{enumerate}
\end{remarks}

\subsection{Immersions régulières et quasi-régulières.}
\EGAReaderTarget{ega4.c04.subsec.16-9}{IV.16.9}

\begin{definition}[16.9.1]
\EGAReaderTarget{ega4.c04.def.16-9-1}{IV.16.9.1}
Soit $X$ un espace annelé. On dit qu'un Idéal $\mathscr{I}$ de
$\mathcal{O}_X$ est \emph{régulier} (resp. \emph{quasi-régulier}) si, pour
tout point $x\in\operatorname{Supp}(\mathcal{O}_X/\mathscr{I})$, il existe un
voisinage ouvert $U$ de $x$ dans $X$ et une suite régulière (0, 15.2.2) (resp.
quasi-régulière (0, 15.2.2)) d'éléments de $\Gamma(U,\mathcal{O}_X)$ qui
engendre $\mathscr{I}|U$.
\end{definition}

On dira qu'une suite régulière (resp. quasi-régulière) de sections de
$\mathcal{O}_X$ au-dessus de $U$ qui engendre $\mathscr{I}|U$ est un
\emph{système régulier} (resp. \emph{quasi-régulier}) \emph{de générateurs}
de $\mathscr{I}|U$.

\begin{definition}[16.9.2]
\EGAReaderTarget{ega4.c04.def.16-9-2}{IV.16.9.2}
Soit $j:Y\to X$ une immersion de préschémas et soit $U$ un ouvert de $X$ tel
que $j(Y)\subset U$ et que $j$ soit une immersion fermée de $Y$ dans $U$. On
dit que $j$ est \emph{régulière} (resp. \emph{quasi-régulière}) si le
sous-préschéma fermé $j(Y)$ de $U$ associé à $j$ est défini par un Idéal
régulier (resp. quasi-régulier) de $\mathcal{O}_U$ (condition indépendante de
l'ouvert $U$ choisi).
\end{definition}

On dit qu'un sous-préschéma $Y$ d'un préschéma $X$ est \emph{régulièrement
immergé} (resp. \emph{quasi-régulièrement immergé}) si l'injection canonique
$j:Y\to X$ est une immersion régulière (resp. quasi-régulière). Si $Y$ est un
sous-préschéma fermé de $X$ et $\mathscr{I}$ l'Idéal de $\mathcal{O}_X$ qui
définit $Y$, il revient au même de dire que $\mathscr{I}$ est \emph{régulier}
(resp. \emph{quasi-régulier}).

Par exemple, si $A$ est un anneau \emph{intègre}, $f$ un élément $\neq0$ de
$A$, le sous-préschéma fermé $V(f)$ de $\operatorname{Spec}(A)$ (isomorphe à
$\operatorname{Spec}(A/fA)$) est \emph{régulièrement immergé} dans
$\operatorname{Spec}(A)$.

Tout Idéal régulier est quasi-régulier (0, 15.2.2); toute immersion régulière
est quasi-régulière (cf. (16.9.11) pour une réciproque).

\begin{proposition}[16.9.3]
\label{IV.16.9.3}
Soient $X$ un espace annelé, $\mathscr{I}$ un Idéal de $\mathcal{O}_X$,
$(f_i)_{1\leq i\leq m}$ une suite finie de sections de $\mathcal{O}_X$
au-dessus de $X$ engendrant $\mathscr{I}$. Pour que $(f_i)$ soit une suite
quasi-régulière (0, 15.2.2), il faut et il suffit que les conditions suivantes
soient vérifiées:
\begin{enumerate}
  \item[(i)] Les images canoniques des $f_i$ dans
  $\mathscr{I}/\mathscr{I}^2$ forment une base de cet
  $\mathcal{O}_X/\mathscr{I}$-Module.
  \item[(ii)] L'homomorphisme surjectif canonique (16.1.2.2)
  \[
    \mathbf{S}_{\mathcal{O}_X/\mathscr{I}}^{\bullet}
    (\mathscr{I}/\mathscr{I}^2)
    \longrightarrow\mathscr{G}r_{\mathscr{I}}^{\bullet}(\mathcal{O}_X)
  \]
  est bijectif.
\end{enumerate}
En outre, s'il en est ainsi, toute suite $(f'_i)_{1\leq i\leq n}$ de $n$
sections de $\mathscr{I}$ au-dessus de $X$ qui engendre $\mathscr{I}$ est
quasi-régulière.
\end{proposition}

